\section{Conclusion}

Despite the rapid advancement of deep learning in hydrological prediction, model interpretability remains a fundamental challenge. Existing explanation approaches are often limited by insufficient verifiability, neglect of variable dependencies\cite{e24050687_Shap_problem}, and a lack of understanding of internal model structures\cite{kumar2020_Shap_problem}. As a result, it remains unclear whether high-performing models truly learn physically meaningful hydrological processes\cite{ReviewOfXAI}.

To address these challenges, this study proposes the Sparse Transformer-based\cite{openai_weight_sparse_transformers} Explainable Hydrology (STEH) framework. The proposed framework establishes an end-to-end interpretability pipeline from inputs, through internal model structures, to outputs. By integrating structured sparsification, circuit-level ablation, variable interaction analysis\cite{wang2022interpretabilitywildcircuitindirect}, and representation probing, STEH provides a systematic approach to uncovering the internal decision-making knowledge structures of deep learning models in hydrological applications.

Overall, the main findings of this study can be summarized as follows.

First, this study provides a circuit-level perspective on deep learning models by explicitly revealing internal decision pathways and information flow. Compared with traditional feature attribution methods, this framework offers a structurally grounded and verifiable paradigm for interpretability.

Second, we identify a counterintuitive yet robust finding: simpler models tend to exhibit more compact internal circuits and stronger physical interpretability, whereas overly complex models often introduce redundant structures that fail to improve—and may even degrade—predictive performance. This result suggests that model parsimony not only benefits generalization but also facilitates the extraction of meaningful physical knowledge.

Third, we demonstrate that physically meaningful representations are consistently learned across models, regardless of predictive performance. Even lower-performing models retain substantial hydrological signals, indicating that physical knowledge is not exclusive to high-accuracy models but is inherently embedded within data-driven models.

Fourth, the STEH framework enables a reliable characterization of variable interactions from a structural perspective. The results reveal hydrologically consistent interaction patterns, including precipitation-dominated processes, temperature-driven dynamics in high-altitude regions, and memory effects associated with antecedent conditions. These interaction structures remain stable across diverse hydro-climatic regimes, suggesting that the learned patterns reflect real-world hydrological knowledge.

Fifth, sensitivity analysis shows that the model exhibits physically consistent responses to precipitation magnitude perturbations, while demonstrating limited sensitivity to temporal redistribution and state perturbations. This indicates that the model effectively captures dominant streamflow-generation knowledge but still lacks a complete representation of temporal organization and state evolution processes.

Sixth, the physical interpretability derived from STEH is further validated through its transferability to a Random Forest model. By incorporating STEH-derived structural information, the modified model achieves consistent and statistically significant performance improvements across stations. This demonstrates that the extracted knowledge is not only physically meaningful but also practically useful.

Nevertheless, several limitations should be acknowledged. First, although STEH improves structural interpretability, its results still cannot be interpreted as strict causal relationships\cite{lipton2018mythos}. This limitation arises from the intrinsic nature of deep learning models as statistical learning systems\cite{pearl2018theoreticalimpedimentsmachinelearning}, where the extracted patterns primarily reflect statistical associations rather than true physical causality. Therefore, the findings derived from the model should be interpreted with caution when relating them to real-world physical processes.

Second, although this study reveals interaction structures among variables, their theoretical linkage to explicit physical process equations remains to be further explored. For example, the specific coupling relationships between variables, the pathways through which interactions operate, and the development of more controllable and interpretable circuit structures are not yet fully characterized\cite{pach2025sparseautoencoderslearnmonosemantic}. Although the current framework enforces sparsity to simplify internal structures, the analysis is still largely limited to the relationship between input variables and the final hidden representations, leaving the information evolution within intermediate high-dimensional feature spaces insufficiently explored.

In addition, the validation of the extracted “physical knowledge” in this study is still based on machine learning models rather than explicit process-based hydrological models\cite{gmd-11-3481-2018}. This may limit the physical rigor of the conclusions, and future work should incorporate process-driven models to further verify the findings.

In summary, this study demonstrates that deep learning models for hydrological prediction do encode physically meaningful information that can be systematically extracted through structural analysis. The proposed STEH framework not only provides a new paradigm for bridging black-box prediction and physical understanding, but also offers a promising pathway for feeding back into the improvement of traditional and physics-based models. These findings suggest that the role of deep learning in hydrology may extend beyond accurate prediction toward serving as a powerful tool for knowledge discovery and advancing the integration of data-driven and process-based modeling.